% Corrigé intégré automatiquement pour la banque Exercices.
% Source : A_Integrer/DDS_01/016_PFS/016_PFS.tex
\begin{corrigebox}[Corrigé]
\CorrigeQuestion{1}
On identifie les classes d'équivalence cinématique, les surfaces de contact et les mouvements relatifs autorisés. Chaque contact est remplacé par la liaison normalisée correspondante, puis les axes et points caractéristiques sont reportés sur le schéma cinématique minimal.
\CorrigeQuestion{2}
La résolution consiste à identifier les données et l'inconnue, choisir la loi du cours adaptée, écrire l'équation symbolique avant toute application numérique, puis vérifier l'homogénéité, le signe et l'ordre de grandeur du résultat. Les hypothèses utilisées doivent être explicitement contrôlées à la fin.
\CorrigeQuestion{3}
On isole le solide ou l'ensemble indiqué, on dresse le bilan complet des actions mécaniques extérieures et on choisit un point de réduction qui élimine le maximum d'inconnues. Le PFS donne $\sum\vec F=\vec0$ et $\sum\vec M_A=\vec0$ ; les projections utiles fournissent les inconnues, puis leur signe permet de vérifier le sens supposé.
\CorrigeQuestion{4}
On écrit la fermeture géométrique sous forme vectorielle, puis on projette dans une base adaptée. Les produits scalaires donnent les relations de longueur et les produits vectoriels les directions normales ; les équations obtenues sont ensuite résolues avec les conditions géométriques du mécanisme.
\CorrigeQuestion{5}
La résolution consiste à identifier les données et l'inconnue, choisir la loi du cours adaptée, écrire l'équation symbolique avant toute application numérique, puis vérifier l'homogénéité, le signe et l'ordre de grandeur du résultat. Les hypothèses utilisées doivent être explicitement contrôlées à la fin.
\CorrigeQuestion{6}
La résolution consiste à identifier les données et l'inconnue, choisir la loi du cours adaptée, écrire l'équation symbolique avant toute application numérique, puis vérifier l'homogénéité, le signe et l'ordre de grandeur du résultat. Les hypothèses utilisées doivent être explicitement contrôlées à la fin.
\CorrigeQuestion{7}
La résolution consiste à identifier les données et l'inconnue, choisir la loi du cours adaptée, écrire l'équation symbolique avant toute application numérique, puis vérifier l'homogénéité, le signe et l'ordre de grandeur du résultat. Les hypothèses utilisées doivent être explicitement contrôlées à la fin.
\end{corrigebox}
